Translating natural language into SQL is nearly solved on curated academic benchmarks, yet the same systems falter on the databases that matter in practice. Three limitations of current approaches explain the gap, and together they define what a system built for industry scale must overcome. First, most LLM pipelines generate a query in a \emph{single pass}, or with only shallow self-correction, so they have no mechanism to recover when a large, unfamiliar schema leads them astray. Second, they implicitly assume the schema is \emph{small enough to fit in context and self-explanatory}, whereas enterprise schemas span thousands of tables and tens of thousands of cryptically named columns that no model can interpret unaided, and that drift and go stale as the business evolves. Third, they are tuned and evaluated against \emph{gold queries}, a luxury that does not exist in production, where the system itself must decide whether a generated query is correct before returning it. Meeting these conditions at industry scale calls for treating schema understanding, retrieval, generation, and validation not as a single forward pass but as fallible steps to be checked and repaired.

This chapter presents our research contribution to date, \emph{Reflect-SQL}~\cite{jain2027reflectsql}, a framework for Text-to-SQL built on exactly this principle and operationalizing the schema-understanding pillar developed in Chapter~\ref{ch:ch1}. As established there, deploying Text-to-SQL in such realistic settings is hindered by three systemic difficulties that mirror the limitations above: (i) \textbf{schema obscurity}, the large-scale databases with cryptic column names (e.g., \texttt{A2}, \texttt{dname}) and undocumented, drifting relationships; (ii) \textbf{contextual retrieval}, the difficulty of extracting the relevant schema elements, from among thousands of candidates, when vague user queries map poorly onto structured schemas; and (iii) \textbf{logical integrity}, the need to move beyond syntactic validity to guarantee that a query faithfully answers the user's intent. Reflect-SQL addresses these through a multi-stage \emph{self-reflection} approach that, rather than trusting a single-pass generation, scores its own intermediate outputs and iteratively refines them within interconnected feedback loops. This directly extends the self-reflection thread of Section~\ref{sec:schema} to the structured domain of databases, exploiting the concrete, checkable signals (executability, syntax, and result inspection) that this domain uniquely affords.

The primary contributions of the framework are: (1) a complete, end-to-end Text-to-SQL system for enterprise settings; (2) a \textbf{dynamic knowledge base} that enriches schema understanding over successful executions; (3) a \textbf{reference-free} methodology for validating SQL without gold-standard labels; and (4) a \textbf{feedback-driven query refinement} process that aligns ambiguous user queries with the database schema. On the challenging BIRD benchmark~\cite{li2023bird}, Reflect-SQL attains an execution accuracy of $72.03\%$, substantially outperforming strong baselines.

\section{Framework Overview}
The framework is organized as a knowledge base feeding three interconnected reflective loops (retrieval, SQL generation, and entailment) as depicted in Fig.~\ref{fig:high-level-arch}. A dynamic \textbf{knowledge base} (KB) acts as a semantic abstraction layer over the raw schema. A \textbf{feedback-driven retrieval} loop then identifies the tables and columns relevant to a query, an \textbf{iterative SQL synthesis} loop generates and corrects a candidate query, and a final \textbf{entailment} check verifies that the executed result matches the user's intent before returning it and enriching the KB.

Each component is shaped by one of the industry-scale pressures identified above, and this is the lens through which the rest of the chapter should be read. The \textbf{knowledge base} answers schema obscurity and context limits: rather than feeding the model a raw schema it cannot fit or interpret, the KB externalizes a semantic layer, descriptions, metadata, and relationships, that the model consults on demand. Its \textbf{iterative refinement} answers schema drift: because an enterprise schema is never fully documented and never stops changing, the KB cannot be built once and frozen, so every validated query feeds its result back to enrich descriptions and relationships, turning deployment into continual learning rather than a one-time setup. The \textbf{SQL synthesis} loop answers the fragility of single-pass generation, catching and correcting syntactic and semantic faults before a query is ever executed. And \textbf{results validation}, through the entailment check, answers the absence of gold labels in production: it judges the executed result against the user's intent in a reference-free manner, so the system can decide for itself whether an answer is trustworthy. We detail each component in turn.

\begin{figure}[ht]
\centering
\resizebox{\textwidth}{!}{%
\begin{tikzpicture}[node distance=1cm and 1.5cm, >=Stealth]
    \colorlet{processFill}{teal!15}
    \colorlet{processBorder}{black!60}
    \colorlet{decisionFill}{orange!20}
    \colorlet{decisionBorder}{black!60}

    \tikzstyle{block} = [rectangle, rounded corners, draw=processBorder, thick, text centered, minimum height=1.5cm, text width=3cm, fill=processFill]
    \tikzstyle{decision} = [diamond, draw=decisionBorder, thick, aspect=1.5, text centered, text width=2.5cm, fill=decisionFill]
    \tikzstyle{arrow} = [->, thick, black!70]
    \tikzstyle{container} = [rectangle, rounded corners, draw, dashed, inner sep=0.5cm, black!80]

    \node[block] (preprocessing) {\small Generate table descriptions based on table schema};
    \node[block, right=of preprocessing] (retrieval) {\small Retrieve relevant tables and columns};
    \node[block, right=of retrieval] (generation) {\small Generate and refine SQL Query};
    \node[decision, right=of generation] (entailment) {Entailed?};
    \node[block, right=of entailment] (postprocessing) {\small Improve Knowledge base with user query and generated SQL};

    \draw [arrow] (preprocessing) -- (retrieval);
    \draw [arrow] (retrieval) -- (generation);
    \draw [arrow] (generation) -- (entailment);
    \draw [arrow] (entailment) -- node[above] {Yes} (postprocessing);
    \draw [arrow] (entailment.north) -- ++(0, 1.2cm) -| node[above, pos=0.25] {No} (retrieval.north);

    \node[container, fit=(preprocessing), label=above:KB Construction] {};
    \node[container, fit=(retrieval), label=above:Retrieval Loop] {};
    \node[container, fit=(generation), label=above:SQL Generation Loop] {};
    \node[container, fit=(postprocessing), label=above:KB Enrichment] {};
\end{tikzpicture}
}
\caption{High-level architecture of the Reflect-SQL framework.}
\label{fig:high-level-arch}
\end{figure}

\section{Knowledge Base Construction and Refinement}
The knowledge base is a semantic abstraction layer containing table descriptions, column metadata, and relationships. It is vital for industry-scale databases, which routinely exceed the context limits of LLMs and suffer from obscure naming conventions. By externalizing schema details into a structured KB, the framework overcomes context constraints and limited schema understanding, ensuring the LLM accesses only the specific information needed to bridge user intent and the database. This design responds directly to the static-KB limitation identified in Section~\ref{sec:schema}: the KB here is not a one-time artifact but is continuously refined.

\subsection{Construction}
For the BIRD dataset~\cite{li2023bird}, the KB is constructed by synthesizing the database schema, sample data values, and the provided training questions and external evidence. In broader enterprise settings, the same process ingests domain knowledge from business documents and glossaries. These inputs are processed to generate comprehensive descriptions for tables, columns, and their relationships. To enable efficient retrieval, the components are encoded into dense vector representations using a sentence embedding model~\cite{reimers2019sentence}. The vector database is structured hierarchically: a global collection for table descriptions and separate local collections for each table's column descriptions. This supports a two-step retrieval process: first identifying relevant tables, then narrowing to specific columns. A representative JSON structure of the KB is shown in Fig.~\ref{fig:kb-json}.

\begin{figure}[ht]
\centering
\begin{minipage}{\linewidth}
\footnotesize
\begin{verbatim}
{
  "financial": {
    "account": {
      "table_description": "it stores information about bank accounts ...",
      "relationships": [{
          "column_name": "district_id", "related_table": "district",
          "related_column": "district_id" }...],
      "column_description": {
        "account_id": "Unique identifier for each bank account...",
        "district_id": "Identifier representing the district...",
        ...}
    }...}
  ...
}
\end{verbatim}
\end{minipage}
\caption{Example JSON specification of the knowledge base.}
\label{fig:kb-json}
\end{figure}

\subsection{Iterative refinement}
In a post-processing step, after a successful SQL execution, one that passes syntactic and semantic validation and whose entailment score exceeds a threshold $\theta_e$, the knowledge base is refined. This conditional update is critical for preventing error propagation: low-quality or contextually incorrect queries could otherwise pollute the KB, so only validated queries with validated results contribute to the system's learning. This mechanism is essential in real-world environments with ambiguous column names and continuously evolving schemas. Table~\ref{tab:kb-refinement} illustrates how column descriptions are enriched after a successful query.

\begin{table}[ht]
\centering
\footnotesize
\begin{tabular}{lp{4.5cm}p{6cm}}
\toprule
\textbf{Component} & \textbf{Before Refinement} & \textbf{After Refinement} \\
\midrule
\texttt{loan.status}    &  Categorical column for loan status with values 'A', 'B', 'C', 'D'.  & Loan repayment state: 'A' indicates finished or paid, 'B' indicates defaulted...  \\
\texttt{trans.k\_symbol} & Text column named 'k\_symbol' with values 'POJISTNE', 'SIPO'. & Transaction Category: determines the payment type (e.g. 'POJISTNE' maps to Insurance)... \\
\bottomrule
\end{tabular}
\caption{Example of knowledge base refinement for BIRD~\cite{li2023bird}.}
\label{tab:kb-refinement}
\end{table}

\section{Iterative RAG Retrieval}
In the retrieval phase, the curated KB is used to identify the tables and columns relevant to a user query, as detailed in Fig.~\ref{fig:retrieval-loop}.

\begin{figure}[ht]
\centering
\resizebox{0.9\columnwidth}{!}{%
\begin{tikzpicture}[node distance=1cm and 0.5cm, >=Stealth]
    \colorlet{processFill}{teal!15}
    \colorlet{processBorder}{black!60}
    \colorlet{decisionFill}{orange!20}
    \colorlet{decisionBorder}{black!60}
    \colorlet{loopBg}{black!5}

    \tikzstyle{io} = [rectangle, draw=processBorder, thick, text centered, minimum height=1cm, text width=2.5cm, fill=processFill]
    \tikzstyle{proc} = [rectangle, rounded corners, draw=processBorder, thick, text centered, minimum height=1cm, text width=3cm, fill=processFill]
    \tikzstyle{decision} = [diamond, draw=decisionBorder, thick, aspect=1.5, text centered, text width=3.5cm, fill=decisionFill]
    \tikzstyle{arrow} = [->, thick, black!70]
    \tikzstyle{group} = [rectangle, rounded corners, draw=processBorder, thick, inner sep=0.3cm, label={[font=\small, black!80]above:Information Retrieval}]

    \node[io] (user_query) {User Query};
    \node[io, above=of user_query] (kb) {Knowledge Base};
    \node[proc, right=1.5cm of user_query] (rag) {RAG Retrieval};

    \node[proc, right=1.5cm of rag, text width=2cm] (tables) {Retrieve tables};
    \node[proc, right=of tables, text width=2cm] (related) {Retrieve related tables};
    \node[proc, right=of related, text width=2cm] (columns) {Retrieve table columns};

    \node[decision, above=1.5cm of related] (scorer) {Retrieved information scorer \\ \texttt{score > threshold}};

    \node[proc, above=1cm of scorer] (prompt) {Prompt to suggest improvements in the user query for better RAG retrieval};
    \node[proc, left=2.5 cm of prompt, text width=2.5cm] (llm) {LLM};
    \node[proc, below=2.5 cm of llm] (improved_query) {Improved user query};

    \node[proc, right=1.5cm of scorer, text width=2.5cm] (sql_gen_loop) {SQL Generation Loop};

    \draw [arrow] (user_query) -- (rag);
    \draw [arrow] (kb) -- (rag);
    \draw [arrow] (rag) -- (tables);
    \draw [arrow] (tables) -- (related);
    \draw [arrow] (related) -- (columns);
    \draw [arrow] (columns.north) -- ++(0,0.5cm) -| (scorer.south);
    \draw [arrow] (scorer.north) -- node[right, pos=0.2] {No} (prompt);
    \draw [arrow] (scorer.east) -- node[above] {Yes} (sql_gen_loop);
    \draw [arrow] (prompt) -- (llm);
    \draw [arrow] (llm) -- (improved_query);
    \draw [arrow] (improved_query) -- (rag);

    \begin{pgfonlayer}{background}
        \node[rectangle, rounded corners, draw=black!80, thick, fill=loopBg, fit=(user_query) (kb) (scorer) (prompt) (columns) (sql_gen_loop), inner sep=0.8cm, label={[font=\bfseries, black!80]above:Retrieval Loop}] {};
        \node[group, fill=white, fit=(tables) (related) (columns)] {};
    \end{pgfonlayer}
\end{tikzpicture}
}
\caption{Iterative retrieval loop.}
\label{fig:retrieval-loop}
\end{figure}

\subsection{Hierarchical table and column retrieval}
Given a user query $Q$, we embed it and use cosine similarity (Eq.~\ref{eq:score_formula}) to retrieve the top-$k$ tables $\{T_1, \dots, T_k\}$ from the vector database.
\begin{equation} \label{eq:score_formula}
\text{score}(T_i, Q) = \cos(\vec{t}_i, \vec{q}) = \frac{\vec{t}_i \cdot \vec{q}}{\|\vec{t}_i\| \|\vec{q}\|}
\end{equation}
This set is expanded by including tables related via relationships defined in the KB, ensuring join paths are available, the join-aware retrieval principle discussed in Section~\ref{sec:schema}. For each selected table, a similar embedding-based search retrieves the most relevant columns.

\subsection{Relevance and completeness scoring}
We assess retrieval quality with two metrics: a \textbf{relevance score} ($R$) for semantic alignment and a \textbf{completeness score} ($C$) to ensure all necessary schema elements are present. These are produced by an LLM-as-a-judge, and the final retrieval score $Score_r$ is a convex combination (Eq.~\ref{eq:retrieval_score}).
\begin{equation} \label{eq:retrieval_score}
Score_r = \alpha R + (1-\alpha) C , \quad \text{where } \quad \alpha \in [0,1]
\end{equation}

\subsection{Query refinement}
If the retrieved schema is insufficient ($Score_r < \theta_r$), an iterative refinement loop is triggered. In each iteration, the LLM reformulates the user query into a more detailed version $Q'$ to sharpen its semantic focus, while the retrieval scope is simultaneously widened by increasing the number of tables and columns retrieved in the next pass. Refining the query's focus while broadening the search area lets the system converge on the complete set of schema elements required for complex queries. The loop terminates once $Score_r \geq \theta_r$ or a maximum number of iterations is reached.

\section{SQL Query Generation}
With a high-quality set of retrieved tables, columns, descriptions, and sample data, a candidate SQL query is generated by an LLM and then subjected to two nested validation loops, as shown in Fig.~\ref{fig:sql-gen-loop}.

\begin{figure}[ht]
\centering
\resizebox{\columnwidth}{!}{%
\begin{tikzpicture}[node distance=1.2cm and 1cm, >=Stealth]
    \colorlet{processFill}{teal!15}
    \colorlet{processBorder}{black!60}
    \colorlet{decisionFill}{orange!20}
    \colorlet{decisionBorder}{black!60}
    \colorlet{loopBg}{black!5}

    \tikzstyle{proc} = [rectangle, rounded corners, draw=processBorder, thick, text centered, minimum height=1cm, text width=4cm, fill=processFill]
    \tikzstyle{io} = [rectangle, draw=processBorder, thick, text centered, minimum height=1cm, text width=2.5cm, fill=processFill]
    \tikzstyle{decision} = [diamond, draw=decisionBorder, thick, aspect=1.5, text centered, text width=3cm, fill=decisionFill]
    \tikzstyle{arrow} = [->, thick, black!70]

    \node[proc] (prompt_initial) {Prompt based on given details (retrieved tables, columns, and relationships) to create SQL for the user query};
    \node[proc, right=1.5cm of prompt_initial, text width=2cm] (llm) {LLM};
    \node[io, below=of llm, text width=2cm] (sql_query) {SQL Query};
    \node[decision, below=of sql_query] (syntax_check) {SQL Syntactically correct?};
    \node[decision, right=1.5cm of syntax_check] (semantic_scorer) {SQL Semantic Scorer \\ \texttt{score > threshold}};
    \node[proc, above=3.2 cm of semantic_scorer, text width=5cm] (prompt_refine) {Prompt LLM to regenerate the accurate query based on semantic scorer feedback};
    \node[proc, left=1 cm of sql_query, text width=5cm] (prompt_refine_error) {Prompt LLM to regenerate the accurate query to correct the error in execution};
    \node[decision, right=1.5cm of semantic_scorer, text width=2.5cm] (entailment_check) {Entailment check};

    \draw [arrow] (prompt_initial) -- (llm);
    \draw [arrow] (llm) -- (sql_query);
    \draw [arrow] (sql_query) -- (syntax_check);
    \draw [arrow] (syntax_check) -- node[above] {Yes} (semantic_scorer);
    \draw [arrow] (prompt_refine_error.east) -- (llm);
    \draw [arrow] (syntax_check.west) -- ++(-2.4, 0) -- node[right, pos=0.8] {No}  (prompt_refine_error.south);
    \draw [arrow] (semantic_scorer.north) -- node[right, pos=0.2] {No} (prompt_refine);
    \draw [arrow] (prompt_refine) -- (llm);
    \draw [arrow] (semantic_scorer.east) -- node[above] {Yes} (entailment_check);

    \begin{pgfonlayer}{background}
        \node[rectangle, rounded corners, draw=black!80, thick, fill=loopBg, fit=(prompt_initial) (prompt_refine) (syntax_check) (semantic_scorer) (entailment_check), inner sep=0.8cm, label={[font=\bfseries, black!80]above:SQL Generation Loop}] {};
    \end{pgfonlayer}
\end{tikzpicture}
}
\caption{Iterative SQL generation loop.}
\label{fig:sql-gen-loop}
\end{figure}

\subsection{Syntactic validation and correction}
Each generated query is immediately checked for syntactic validity using a database engine. On a syntax error, a correction loop feeds the original query, the error message, and the schema context back to the LLM with instructions to fix the error. The loop repeats until a syntactically correct query is produced or a maximum iteration limit is met, making the system resilient to basic generation mistakes.

\subsection{Semantic validation}
A syntactically valid query undergoes semantic validation across four dimensions: column-name consistency ($C$), granularity and aggregation ($G$), join and relationship consistency ($J$), and query-logic alignment ($Q$). An LLM-as-a-judge scores each dimension, and a final semantic score $S$ is a weighted sum (Eq.~\ref{eq:semantic_formula}).
\begin{equation} \label{eq:semantic_formula}
S = w_c \cdot C + w_g \cdot G + w_j \cdot J + w_q \cdot Q
\end{equation}
where $w_c, w_g, w_j, w_q$ are predefined weights. If $S < \theta_s$, the query is returned to the LLM with diagnostic feedback for iterative correction.

\section{Entailment}
Once a query passes all validation, it is executed to retrieve the result set $R$. In a final verification stage, the framework checks whether $R$ logically entails the user's original query $Q$, the reference-free validation that distinguishes this pillar's notion of \emph{logical integrity} from mere syntactic correctness.

\subsection{Entailment scoring}
An entailment score $E(Q, R)$ is computed by an LLM-as-a-judge, assessing both \textbf{schema consistency} ($E_{\text{schema}}$) and \textbf{data consistency} ($E_{\text{data}}$), as in Eq.~\ref{eq:entailment_score}.
\begin{equation} \label{eq:entailment_score}
E(Q, R) = \beta \cdot E_{\text{schema}} + (1 - \beta) \cdot E_{\text{data}} \quad \text{where } \quad \beta \in [0,1]
\end{equation}

\subsection{Entailment decision}
If $E(Q, R) \geq \theta_e$, the result is returned to the user and the KB is updated with the successful query--result pair. Otherwise, the entire pipeline is re-triggered with feedback on the failure reason (e.g., missing columns, incorrect filters). This final loop ensures the output truly matches the user's intent.

\section{Mitigating Bias in LLM-based Evaluation}
Using an LLM as a judge can introduce biases~\cite{zheng2023judging}. To mitigate this, the framework employs a structured, few-shot prompting strategy for all scoring tasks: prompts provide detailed instructions, scoring rubrics, and few-shot examples that anchor the LLM's judgment to task-aligned criteria, increasing the reliability and reproducibility of the evaluation while reducing the influence of the model's inherent biases.

\section{Model Parameter Optimization}
\label{sec:parameters}
A systematic grid search determined the optimal thresholds ($\theta_r, \theta_s, \theta_e$) and weights ($\alpha, \beta, w_c, w_g, w_j, w_q$). Each parameter was varied and the combination yielding the highest validation accuracy was selected; the final values are given in Table~\ref{tab:optimized-params}.

\begin{table}[ht]
\centering
\begin{tabular}{ll}
\toprule
\textbf{Parameter} & \textbf{Value} \\
\midrule
$\theta_r$ & 4 \\
$\theta_s$ & 5 \\
$\theta_e$ & 3 \\
$\alpha$   & 0.6 \\
$\beta$    & 0.2 \\
$w_c$      & 0.1 \\
$w_g$      & 0.3 \\
$w_j$      & 0.4 \\
$w_q$      & 0.2 \\
\bottomrule
\end{tabular}
\caption{Optimized parameters using grid search.}
\label{tab:optimized-params}
\end{table}

\section{Evaluation}
\subsection{Experimental setup}
We evaluate the framework on the \textbf{BIRD benchmark}~\cite{li2023bird}. Because large-scale, proprietary enterprise datasets are generally inaccessible for academic experiments, BIRD serves as an excellent proxy: it tests Text-to-SQL systems over diverse relational databases that mirror real-world complexity.

\subsection{Results}
We evaluate Claude, OpenAI, and DeepSeek API-based models, as well as an open-source DeepSeek model. Table~\ref{tab:eval-results} reports execution accuracy on BIRD. Across every model, enabling the reflection framework produces substantial gains over the corresponding non-reflective baseline (a 6--8\% overall improvement), confirming that iterative reasoning and refinement are central to performance.

\begin{table}[ht]
\centering
\resizebox{\columnwidth}{!}{%
\begin{tabular}{lcccc}
\toprule
\textbf{Model} & \textbf{Simple} & \textbf{Moderate} & \textbf{Challenging} & \textbf{Overall} \\
\midrule
\multicolumn{5}{l}{\textbf{Claude Models}} \\
\textbf{Claude-sonnet-4.5 (with reflection)}  & \textbf{88.97} & \textbf{50.65} & \textbf{32.41} & \textbf{72.03} \\
\texttt{Claude-sonnet-4.5 (w/o reflection)}     & 81.18 & 45.25  & 24.13 & 64.92 \\
\midrule
\multicolumn{5}{l}{\textbf{OpenAI API Models}} \\
\textbf{o1-preview (with reflection)}  & \textbf{75.89} & \textbf{48.49} & \textbf{18.62} & \textbf{62.19} \\
\texttt{o1-preview (w/o reflection)}     & 69.35 & 40.54  & 15.86 & 55.66 \\
\midrule
\multicolumn{5}{l}{\textbf{DeepSeek API Models}} \\
\textbf{deepseek-reasoner (with reflection)}   & \textbf{73.62} & \textbf{46.55} & \textbf{17.24} & \textbf{60.10} \\
\texttt{deepseek-reasoner} (w/o reflection)  & 62.90 & 40.54  & 12.41 & 51.88 \\
\midrule
\multicolumn{5}{l}{\textbf{Open Source Models (Local setup)}} \\
\textbf{DeepSeek-Coder-V2-Lite... (with reflection)}    & \textbf{43.55} & \textbf{21.62} & \textbf{3.44} & \textbf{35.02} \\
\texttt{DeepSeek-Coder-V2-Lite...} (w/o reflection)    & 30.65 & 13.51 & 0.0 & 22.64 \\
\bottomrule
\end{tabular}
}
\caption{Accuracy (\%) of various models on the BIRD dataset. Top performers in each category are highlighted.}
\label{tab:eval-results}
\end{table}

\subsection{Human evaluation}
To complement the automatic evaluation, we conducted a qualitative human assessment of the correctness and semantic alignment of generated queries, focusing on the impact of the entailment check on subtle logical errors that syntax or basic semantic validation may miss. Table~\ref{tab:refined-sql} shows a representative case in which entailment feedback corrects a division-by-zero and casting error that would otherwise return null results.

\begin{table}[ht]
\centering
\small
\begin{tabular}{ >{\bfseries}m{2.2cm} m{9.5cm} }
\toprule
\textbf{Stage} & \textbf{Analysis / SQL Snippet} \\
\midrule
User Query & List the lowest three eligible free rates for students in continuation schools. \\
\addlinespace
Initial SQL & \texttt{SELECT ... (f.Count / f.Enrollment) ... WHERE s.SOCType LIKE '\%Continuation\%'} \\
\addlinespace
Entailment Check Result & \textbf{Score}: 2 / 5 \newline \textbf{Reason:} The result returns None values due to division by zero or nulls. Filters for positive enrollment and float casting are needed. \\
\addlinespace
Refined SQL & \texttt{SELECT ... (CAST(f.Count AS REAL) / f.Enrollment) ... WHERE f.Enrollment > 0 AND s.SOCType LIKE '\%Continuation\%'} \\
\bottomrule
\end{tabular}
\caption{Example of SQL refinement driven by entailment feedback.}
\label{tab:refined-sql}
\end{table}

\subsection{Comparison with state of the art}
We benchmark against leading methods on the BIRD development set (Table~\ref{tab:sota}). Reflect-SQL achieves an execution accuracy of $72.03\%$, a marked improvement over DIN-SQL~\cite{pourreza2023din} ($50.72\%$), DAIL-SQL~\cite{gao2023text} ($57.41\%$), and MAC-SQL~\cite{wang2024mac} ($59.59\%$), and is comparable to the recent CHASE-SQL architecture~\cite{pourreza2024chase}, establishing Reflect-SQL as a top-tier solution for complex database querying.

\begin{table}[ht]
\centering
\begin{tabular}{lc}
\toprule
\textbf{Method} & \textbf{Overall Accuracy} \\
\midrule
BIRD paper~\cite{li2023bird} & 46.35 \\
DIN-SQL~\cite{pourreza2023din} & 50.72 \\
DAIL-SQL~\cite{gao2023text} & 57.41 \\
MAC-SQL~\cite{wang2024mac} & 59.59 \\
CHASE-SQL + Claude 3.5 Sonnet~\cite{pourreza2024chase} & 69.53 \\
CHASE-SQL + Gemini 1.5 Pro~\cite{pourreza2024chase} & \textbf{73.01} \\
\midrule
\textbf{Reflect-SQL + Claude 4.5 Sonnet (Ours)} & \textbf{72.03} \\
\bottomrule
\end{tabular}
\caption{Comparison with state-of-the-art methods on the BIRD dev set (execution accuracy \%).}
\label{tab:sota}
\end{table}

\section{Ablation Study}
To quantify each component's contribution, we performed an ablation using \texttt{deepseek-chat} as the base model (Table~\ref{tab:ablation}). Removing any component degrades performance, and the self-reflection loop in SQL generation proves most critical: its removal causes the sharpest drop. These results affirm the effectiveness of the multi-stage, reflective design.

\begin{table}[ht]
\centering
\footnotesize
\resizebox{\columnwidth}{!}{%
\begin{tabular}{lcccc}
\toprule
\textbf{Variant} & \textbf{Simple} & \textbf{Moderate} & \textbf{Challenging} & \textbf{Overall} \\
\midrule
\textbf{Full Model (Self-Reflect)}       & 68.64 & 43.10 & 14.48 & 55.80 \\
w/o Self-Reflection in SQL Generation    & 62.91 & 37.93 & 12.41 & 50.58 \\
w/o Iterative Retrieval                  & 67.02 & 40.30 & 10.34 & 53.58 \\
w/o Entailment Check                     & 64.10 & 41.59 & 13.10 & 52.47 \\
\bottomrule
\end{tabular}
}
\caption{Ablation results showing the contribution of each module.}
\label{tab:ablation}
\end{table}

\section{Summary}
Reflect-SQL demonstrates that self-reflection, adapted to the structured domain of databases, is an effective route to reliable schema understanding. Its dynamic knowledge base disambiguates obscure schemas and adapts over time; its feedback-driven retrieval aligns vague queries with the schema; and its layered generation and entailment loops enforce logical integrity without gold-standard labels. The framework's effectiveness depends on the initial quality of the knowledge base and can face scalability challenges with extremely large schemas, and the automatic metrics for semantic correctness and entailment remain an evolving area of research. These observations, together with the data-understanding challenges of Chapter~\ref{ch:ch1}, motivate the research direction discussed in the concluding chapter.
